\documentclass{amsart}

\usepackage{amssymb}
\usepackage{booktabs}
\usepackage[hyphens]{url}
\usepackage[pdfauthor={Knowledge Sense},
  pdftitle={A uniform discrepancy bound for anti-recurrence sequences},
  pdfsubject={Mathematics},
  pdfkeywords={anti-recurrence sequence, complementary sequences, mex, discrepancy, automatic sequence},
  colorlinks=false,hidelinks]{hyperref}

\hyphenation{manuscript}

\newtheorem{theorem}{Theorem}[section]
\newtheorem{lemma}[theorem]{Lemma}
\newtheorem{proposition}[theorem]{Proposition}
\newtheorem{corollary}[theorem]{Corollary}
\newtheorem{conjecture}[theorem]{Conjecture}
\theoremstyle{definition}
\newtheorem{definition}[theorem]{Definition}
\theoremstyle{remark}
\newtheorem{remark}[theorem]{Remark}

\numberwithin{equation}{section}
\numberwithin{table}{section}

\newcommand{\N}{\mathbb{N}}
\newcommand{\ba}{\mathbf{a}}

\begin{document}

\title[Discrepancy of anti-recurrence sequences]{A uniform discrepancy bound\\ for anti-recurrence sequences}

\author{Knowledge Sense}
\address{Knowledge Sense}
\email{math@knowledgesense.jp}

\date{13 October 2026}

\subjclass[2020]{Primary 11B83; Secondary 11B37, 11B85}

\keywords{Anti-recurrence sequence, complementary sequences, mex, discrepancy, automatic sequence}

\begin{abstract}
Let $k\ge 2$ and let $a_1,\dots,a_k$ be positive integers with sum $\tau$. The anti-recurrence sequence generated by $(a_1,\dots,a_k)$ is the unique strictly increasing sequence $(A_n)_{n\ge1}$ of positive integers whose complement $(B_n)_{n\ge1}$ in the positive integers satisfies $A_n=\sum_{j=1}^k a_jB_{k(n-1)+j}$ for all $n\ge1$. Fokkink and Joshi conjectured that $A_n-\kappa n$, where $\kappa=k\tau+1$, is a $\tau$-automatic sequence, and proved this under an additional upper bound on $A_1$. We prove that $\lvert A_n-\kappa n\rvert\le 2\kappa$ for every $n\ge1$ and every positive coefficient vector. This is the weak form of the conjecture in which automaticity is replaced by boundedness. The proof combines the classical rank bracketing for complementary sequences with an exact identity for the weighted block sum, and leads to an estimate in which the discrepancy at index $n$ is controlled, with weight $1/k$, by discrepancies at strictly earlier indices.
\end{abstract}

\maketitle

\section{Introduction}

Two strictly increasing sequences of positive integers are \emph{complementary} if every positive integer is a term of exactly one of them. Following Fokkink and Joshi~\cite{FJ}, let $k\ge2$ and let $\ba=(a_1,\dots,a_k)$ be a vector of positive integers, with \emph{trace} $\tau=a_1+\dots+a_k$. A pair $A=(A_n)_{n\ge1}$, $B=(B_n)_{n\ge1}$ of complementary sequences such that
\begin{equation}\label{eq:def}
A_n=\sum_{j=1}^{k}a_jB_{k(n-1)+j}\qquad(n\ge1)
\end{equation}
exists and is unique~\cite[Lemma~2]{FJ}; $A$ is the \emph{anti-recurrence sequence} generated by $\ba$. Concretely, the $n$-th block $B_{k(n-1)+1},\dots,B_{kn}$ consists of the $k$ smallest positive integers not yet used, and $A_n$ is their weighted sum; in other words, each term of the block is the mex (minimal excluded value), that is, the smallest positive integer not among the numbers already used. For $\ba=(1,1)$ one obtains the anti-Fibonacci numbers $3,9,13,18,23,\dots$ (A075326 in~\cite{OEIS}); for $\ba=(1,2)$ and $\ba=(2,1)$ the anti-Pell and anti-Jacobsthal numbers (A304502 and A304499); and for $\ba=(1,1,1)$ and $\ba=(1,1,1,1)$ the anti-Tribonacci and anti-Tetrabonacci numbers (A265389 and A299409). Kimberling and Moses~\cite[Section~6]{KM} described the anti-Fibonacci numbers and their generalizations through systems of complementary equations, and conjectured morphic descriptions of the associated difference sequences.

Put $\kappa=k\tau+1$ and define the \emph{discrepancy} of $A$ by
\[
D_n=A_n-\kappa n\qquad(n\ge1).
\]
Fokkink and Joshi formulated the following conjecture~\cite[Conjecture~2]{FJ}.

\begin{conjecture}[Fokkink--Joshi]\label{conj:FJ}
For every positive integral vector $\ba$ of dimension $k\ge2$, the sequence $(D_n)$ is $\tau$-automatic.
\end{conjecture}

An automatic sequence takes only finitely many values, so Conjecture~\ref{conj:FJ} implies that $(D_n)$ is bounded. Fokkink and Joshi call this boundedness statement a weak form of the conjecture, and note that this is how Kimberling and Moses formulate their conjectures~\cite[p.~6]{FJ}. The conjecture has been proved in the following cases. For $\ba=(1,1)$, Zaslavsky~\cite[Theorem~1]{Z} gave an explicit formula; in the present indexing it gives $A_n-5n\in\{-2,-1\}$. Bosma et al.~\cite[Section~4]{B} treated $\ba=(1,1,1)$ and $\ba=(1,1,1,1)$ with the automatic theorem prover \texttt{Walnut} (see \cite[Theorem~12]{B} for the first). Fokkink and Joshi proved the conjecture for $\ba=(1,2)$ and $\ba=(2,1)$ \cite[Theorems~1 and~2]{FJ}, for the anti-bonacci vectors $\ba=(1,\dots,1)$ \cite[Section~4 and Theorem~3]{FJ}, and for every vector satisfying
\begin{equation}\label{eq:A1bounded}
A_1\le(k-1)\tau+2,
\end{equation}
which they call \emph{$A_1$-bounded} \cite[Definition~1 and Theorem~4]{FJ}. All the special vectors above satisfy~\eqref{eq:A1bounded}. The proof of the last result rests on the fact that, under~\eqref{eq:A1bounded}, every $A_n$ lies in the interval $[\kappa(n-1)+1,\kappa n]$ \cite[Lemma~5]{FJ}. For $\ba=(1,3)$, which does not satisfy~\eqref{eq:A1bounded}, Fokkink and Joshi display a guessed $4$-DFAO (an automaton reading base-$4$ expansions) for the difference sequence and describe how the \texttt{Walnut} verification used for $(1,2)$ and $(2,1)$ can be applied to it \cite[Section~5]{FJ}; this is not stated there as a theorem. They ask whether the conjecture holds without the restriction~\eqref{eq:A1bounded} \cite[Section~6]{FJ}.

Our main result establishes the weak form of Conjecture~\ref{conj:FJ} for all positive coefficient vectors.

\begin{theorem}\label{thm:main}
Let $k\ge2$, let $a_1,\dots,a_k$ be positive integers, and put $\tau=\sum_{j}a_j$ and $\kappa=k\tau+1$. Then the anti-recurrence sequence generated by $(a_1,\dots,a_k)$ satisfies
\[
\lvert A_n-\kappa n\rvert\le 2\kappa\qquad\text{for every }n\ge1.
\]
\end{theorem}

No condition is imposed on the order or the size of the coefficients, and the bound depends only on $k$ and $\tau$. We do not claim that the constant $2\kappa$ is sharp; the ranges observed numerically are considerably smaller (Section~\ref{sec:numerics}). Theorem~\ref{thm:main} does not address automaticity; it reduces Conjecture~\ref{conj:FJ} to a statement about a sequence over a finite alphabet.

Without~\eqref{eq:A1bounded}, the fixed-interval property of~\cite[Lemma~5]{FJ} fails in general. For $\ba=(1,4)$ we have $\tau=5$, $\kappa=11$ and $A_1=9>7=(k-1)\tau+2$, and direct computation gives $A_{348}=3829=11\cdot348+1$, so $A_{348}\notin[\kappa\cdot 347+1,\kappa\cdot348]$. This is consistent with the observation in~\cite[Section~5]{FJ} that $A_{5n+1}=9+55n$ holds up to $n=347$ but not for $n=348$; indeed $A_{1741}=19148\ne19149$. More generally, the vectors $\ba=(1,d)$ with $d\ge3$, whose anti-recurrences Fokkink and Joshi call rusty numbers, are not $A_1$-bounded, since $A_1=2d+1>d+3$. Theorem~\ref{thm:main} shows that their discrepancy sequences are bounded as well.

\emph{Method.} Let $b=B_{k(n-1)+1}$ be the first term of the $n$-th block and let $h$ be the number of terms of $A$ below $b$. The bracketing $A_h<b<A_{h+1}$, together with $b=k(n-1)+1+h$, is the classical rank inequality for complementary sequences; see Fraenkel and Peled~\cite[Theorem~3.2]{FP}. For complementary sequences $(a_n)_{n\ge0}$ and $(b_n)_{n\ge0}$ in their setting, Fraenkel and Peled also show that both sequences are approximately linear if and only if the same-index difference $b_n-a_n$ is approximately linear~\cite[Theorem~4.3]{FP}. For anti-recurrences no such information on a same-index difference is available a priori, since $A_n$ is a weighted sum of a block of $k$ terms of $B$, and we argue differently. The new ingredient is the exact identity
\[
A_n=k\tau n+\tau h+A_1-k\tau+\delta_n,\qquad 0\le\delta_n\le\tau,
\]
in which $\delta_n$ is a block correction (Lemma~\ref{lem:block}). Combined with the bracketing, it yields the estimate
\[
\lvert D_n\rvert\le K+\frac1k\max_{0\le i<n}\lvert D_i\rvert\qquad(n\ge3)
\]
for a constant $0<K\le\kappa$ (Proposition~\ref{prop:envelope}). For $n\ge3$ the indices $h$ and $h+1$ at which the bracketing is applied are strictly smaller than $n$ (Lemma~\ref{lem:earlier}), so the estimate can be used inductively, and Theorem~\ref{thm:main} follows.

\emph{Dependencies.} The basic facts we use, namely the existence and uniqueness of anti-recurrences~\cite[Lemma~2]{FJ}, the inequality of~\cite[Lemma~1]{FJ} and the rank bracketing of~\cite[Theorem~3.2]{FP}, are reproved in Section~\ref{sec:prelim} in the form in which we need them.

\emph{Organization.} Section~\ref{sec:prelim} collects the preliminary facts. Section~\ref{sec:block} proves the block identity and the earlier-index property. Section~\ref{sec:proof} proves Theorem~\ref{thm:main}. Section~\ref{sec:numerics} reports finite numerical checks.

\section{Preliminaries}\label{sec:prelim}

Throughout, $k\ge2$ and $\ba=(a_1,\dots,a_k)$ is a vector of positive integers with trace $\tau$. We identify a strictly increasing sequence with the set of its terms when convenient. For $n\ge0$ write
\[
S_n=\{A_1,\dots,A_n\}\cup\{B_1,\dots,B_{kn}\},
\]
so that $S_0=\emptyset$.

\begin{lemma}[{\cite[Lemma~2]{FJ}}]\label{lem:exist}
There is exactly one pair of complementary sequences $A$, $B$ satisfying~\eqref{eq:def}. Moreover, for every $n\ge1$:
\begin{enumerate}
\item[(i)] $A_n>B_{kn}$;
\item[(ii)] $B_{k(n-1)+1},\dots,B_{kn}$ are the $k$ smallest positive integers not in $S_{n-1}$.
\end{enumerate}
\end{lemma}

\begin{proof}
We first show that every solution satisfies (i) and (ii). Since $a_k\ge1$ and $k\ge2$,
\[
A_n\ge a_kB_{kn}+\sum_{j=1}^{k-1}a_jB_{k(n-1)+j}>B_{kn}.
\]
Let $x\le B_{kn}$ be a positive integer not in $S_{n-1}$. If $x$ were a term $A_i$, then $A_i\le B_{kn}<A_n$ would give $i\le n-1$ and $x\in S_{n-1}$. Hence $x=B_j$ with $k(n-1)<j\le kn$. Thus the integers in $[1,B_{kn}]$ outside $S_{n-1}$ are exactly $B_{k(n-1)+1},\dots,B_{kn}$, which proves (ii). By (ii) and~\eqref{eq:def}, the set $S_{n-1}$ determines the $n$-th block and $A_n$, and hence $S_n$; by induction on $n$ the solution is unique.

For existence, define the blocks recursively by (ii) and the terms $A_n$ by~\eqref{eq:def}, and put $B_0=0$. Assume that $[1,B_{k(n-1)}]\subseteq S_{n-1}$, which is trivial for $n=1$. Then the $n$-th block lies above $B_{k(n-1)}$, and since it consists of the smallest integers outside $S_{n-1}$, we get $[1,B_{kn}]\subseteq S_{n-1}\cup\{B_{k(n-1)+1},\dots,B_{kn}\}$. The sequence $B$ constructed so far is strictly increasing, so $B_{k(n-1)+j}-B_{k(n-2)+j}\ge k$ for $n\ge2$ and $1\le j\le k$, and therefore $A_n-A_{n-1}\ge k\tau>0$. Together with $A_n>B_{kn}$, proved as above, this shows that $A_n$ exceeds every element of $S_{n-1}$ and every term of the $n$-th block. Hence $A_n$ is new, $[1,B_{kn}]\subseteq S_n$, and the induction continues. Terms of later blocks are chosen outside the sets $S_n$, so $A$ and $B$ are disjoint and strictly increasing, and since $B_{kn}\ge kn$, every positive integer lies in some $S_n$.
\end{proof}

\begin{lemma}\label{lem:initial}
The following hold.
\begin{enumerate}
\item[(a)] $B_j=j$ for $1\le j\le k$ and $A_1=\sum_{j=1}^k ja_j$. In particular $\tau\le A_1\le k\tau$, $A_1\ge k(k+1)/2\ge k+1$ and $\tau\ge k$.
\item[(b)] {\rm(\cite[Lemma~1]{FJ})} $A_{n+1}-A_n\ge k\tau\ge k^2>k$ for every $n\ge1$.
\end{enumerate}
Consequently, every set of $k+1$ consecutive integers contains at most one term of $A$.
\end{lemma}

\begin{proof}
(a) By Lemma~\ref{lem:exist}(ii) with $n=1$, the first block is $\{1,\dots,k\}$. The bounds follow from $1\le j\le k$ and $a_j\ge1$, and $k(k+1)/2\ge k+1$ because $k\ge2$.

(b) Since $B$ is a strictly increasing sequence of integers, $B_{kn+j}-B_{k(n-1)+j}\ge k$ for $1\le j\le k$; multiplying by $a_j$ and summing gives $A_{n+1}-A_n\ge k\tau$. Finally $k\tau\ge k^2>k$ by~(a).
\end{proof}

It is convenient to put $A_0=0$. This is only a notational convention, and $0$ is not a term of either sequence.

\begin{lemma}[Rank bracketing]\label{lem:rank}
Let $m\ge1$, let $b=B_m$ and let $h=\#\{i\ge1:A_i<b\}$. Then
\[
b=m+h\qquad\text{and}\qquad A_h<b<A_{h+1}.
\]
\end{lemma}

\begin{proof}
The integers $1,\dots,b-1$ consist of the $m-1$ terms $B_1,\dots,B_{m-1}$ and the $h$ terms of $A$ below $b$, so $b-1=(m-1)+h$. By the definition of $h$ we have $A_h<b\le A_{h+1}$ (for $h=0$ the first inequality reads $0<b$), and $A_{h+1}\ne b$ because $b\in B$.
\end{proof}

Lemma~\ref{lem:rank} is the second inequality $a_{b_k-k}<b_k<a_{b_k-k+1}$ of~\cite[Theorem~3.2]{FP}, transcribed to the present indexing; in~\cite{FP} both sequences start with the term $0$, which corresponds to our convention $A_0=0$.

\section{The block identity}\label{sec:block}

Fix $n\ge1$ and put
\[
m=k(n-1)+1,\qquad b=B_m,\qquad h=\#\{i\ge1:A_i<b\}.
\]
Thus $b$ is the first term of the $n$-th block, and $h=h_n$ depends on $n$. For $1\le j\le k$ define
\[
\varepsilon_j=B_{m+j-1}-b-(j-1),\qquad \delta_n=\sum_{j=1}^{k}a_j\varepsilon_j.
\]

\begin{lemma}\label{lem:block}
With this notation:
\begin{enumerate}
\item[(a)] $b=k(n-1)+1+h$ and $0\le h\le n-1$;
\item[(b)] $\varepsilon_1=0$ and $\varepsilon_j\in\{0,1\}$ for $1\le j\le k$; in particular $0\le\delta_n\le\tau$;
\item[(c)] $A_n=k\tau n+\tau h+A_1-k\tau+\delta_n$.
\end{enumerate}
\end{lemma}

\begin{proof}
(a) The identity is Lemma~\ref{lem:rank}. Since $b\le B_{kn}<A_n$ by Lemma~\ref{lem:exist}(i), no term $A_i$ with $i\ge n$ lies below $b$, so $h\le n-1$.

(b) By Lemma~\ref{lem:initial}, the interval $[b,b+k]$ contains at most one term of $A$, and hence at least $k$ terms of $B$. These are terms $B_i$ with $i\ge m$, since $b=B_m$; therefore $B_m,\dots,B_{m+k-1}$ all lie in $[b,b+k]$. For $1\le j\le k$, the interval $[b,B_{m+j-1}]$ contains exactly $j$ terms of $B$, so $\varepsilon_j$ is the number of terms of $A$ in this interval. As the interval is contained in $[b,b+k]$, we get $\varepsilon_j\in\{0,1\}$, and $\varepsilon_1=0$ by definition. Hence $0\le\delta_n\le\tau$.

(c) By~\eqref{eq:def} and Lemma~\ref{lem:initial}(a),
\[
A_n=\sum_{j=1}^k a_j\bigl(b+j-1+\varepsilon_j\bigr)=\tau b+(A_1-\tau)+\delta_n.
\]
Substituting $b=k(n-1)+1+h$ from (a) gives the claim.
\end{proof}

The next lemma shows that, from the third block on, both indices in the bracketing $A_h<b<A_{h+1}$ are smaller than $n$.

\begin{lemma}\label{lem:earlier}
If $n\ge3$, then $h\le n-2$, that is, $h+1\le n-1$.
\end{lemma}

\begin{proof}
By Lemma~\ref{lem:initial}, $A_1\ge k+1$ and $A_{i+1}-A_i\ge k^2\ge2k$, so
\[
A_{n-1}\ge k+1+2k(n-2).
\]
By Lemma~\ref{lem:block}(a), $b\le k(n-1)+1+(n-1)=(k+1)n-k$. Hence
\[
A_{n-1}-b\ge(k-1)n-2k+1\ge(k-1)\cdot3-2k+1=k-2\ge0.
\]
Since $b$ is a term of $B$, it is different from $A_{n-1}$, so $A_{n-1}>b$. Thus at most $n-2$ terms of $A$ lie below $b$.
\end{proof}

\section{Proof of Theorem~\ref{thm:main}}\label{sec:proof}

Recall that $D_n=A_n-\kappa n$, and put $D_0=0$, in accordance with $A_0=0$. Let
\[
c=A_1-k\tau,\qquad K=-c+\tau+1=(k+1)\tau-A_1+1.
\]
By Lemma~\ref{lem:initial}(a) we have $c\le0$ and $\tau+1\le K\le\kappa$; in particular $0<K\le\kappa$. For $n\ge1$ let
\[
M_{n-1}=\max_{0\le i\le n-1}\lvert D_i\rvert.
\]
This is a maximum over finitely many indices.

\begin{proposition}\label{prop:envelope}
For every $n\ge1$, with $h=h_n$ as in Section~\ref{sec:block},
\begin{equation}\label{eq:two-sided}
c-\tau-1+\frac{1-D_{h+1}}{k}\;\le\;D_n\;\le\;c+\tau-1-\frac{D_h}{k}.
\end{equation}
Consequently,
\begin{equation}\label{eq:envelope}
\lvert D_n\rvert\le K+\frac{M_{n-1}}{k}\qquad(n\ge3).
\end{equation}
\end{proposition}

\begin{proof}
Since $\kappa n=k\tau n+n$, Lemma~\ref{lem:block}(c) can be written as
\begin{equation}\label{eq:Dn}
D_n=\tau h-n+c+\delta_n.
\end{equation}
By Lemma~\ref{lem:rank} applied to $b=B_m$, we have $A_h\le b-1$ and $A_{h+1}\ge b+1$; the first inequality holds also for $h=0$, since $A_0=0$ and $b\ge1$. We write $A_i=\kappa i+D_i$ and use $b=k(n-1)+1+h$. The first inequality becomes $(k\tau+1)h+D_h\le k(n-1)+h$, that is,
\[
\tau h-n\le-1-\frac{D_h}{k}.
\]
The second becomes $(k\tau+1)(h+1)+D_{h+1}\ge k(n-1)+h+2$, that is,
\[
\tau h-n\ge-\tau-1+\frac{1-D_{h+1}}{k}.
\]
Inserting these two bounds into~\eqref{eq:Dn} and using $0\le\delta_n\le\tau$ from Lemma~\ref{lem:block}(b) gives~\eqref{eq:two-sided}.

Now let $n\ge3$. By Lemma~\ref{lem:earlier}, both $h$ and $h+1$ lie in $\{0,\dots,n-1\}$, so $\lvert D_h\rvert\le M_{n-1}$ and $\lvert D_{h+1}\rvert\le M_{n-1}$. Since $c\le0$, we have $c+\tau-1\le-c+\tau+1=K$, and the right-hand inequality in~\eqref{eq:two-sided} gives $D_n\le K+M_{n-1}/k$. The left-hand inequality gives
\[
D_n\ge c-\tau-1+\frac1k-\frac{M_{n-1}}{k}>-K-\frac{M_{n-1}}{k}.
\]
This proves~\eqref{eq:envelope}.
\end{proof}

\begin{proof}[Proof of Theorem~\ref{thm:main}]
We show $\lvert D_n\rvert\le2\kappa$ by strong induction on $n$.

For $n=1$, Lemma~\ref{lem:initial}(a) gives $\tau-\kappa\le D_1=A_1-\kappa\le k\tau-\kappa=-1$, so $\lvert D_1\rvert<\kappa$.

For $n=2$, Lemma~\ref{lem:block}(a) gives $h_2\le1$ and hence $b=k+1+h_2\le k+2$. By Lemma~\ref{lem:block}(b), every term of the second block is at most $b+k\le2k+2$, so $A_2\le(2k+2)\tau$ by~\eqref{eq:def}. Therefore
\[
1-2\kappa\le D_2\le(2k+2)\tau-2\kappa=2\tau-2,
\]
and $\lvert D_2\rvert\le2\kappa$.

Let $n\ge3$ and assume $\lvert D_i\rvert\le2\kappa$ for $1\le i\le n-1$. Since also $D_0=0$, we have $M_{n-1}\le2\kappa$, and~\eqref{eq:envelope} together with $K\le\kappa$ and $k\ge2$ gives
\[
\lvert D_n\rvert\le\kappa+\frac{2\kappa}{k}\le2\kappa.
\]
This completes the induction.
\end{proof}

\begin{remark}
Lemma~\ref{lem:earlier} is what makes~\eqref{eq:envelope} usable for induction. By Lemma~\ref{lem:block}(a) we always have $h+1\le n$, and for $n\le2$ the case $h+1=n$ can occur; the left-hand bound in~\eqref{eq:two-sided} then involves $D_n$ itself. For this reason the cases $n=1,2$ are treated directly.
\end{remark}

\section{Numerical checks}\label{sec:numerics}

The computations reported in this section are not used in the proof. We generated anti-recurrence sequences directly from the mex description of Lemma~\ref{lem:exist}, using exact integer arithmetic. For every generated index $n$ we verified Lemma~\ref{lem:exist}(i), Lemma~\ref{lem:block} and the bound of Theorem~\ref{thm:main}, and for every pair of consecutive generated terms the inequality of Lemma~\ref{lem:initial}(b); for every generated index $n\ge3$ we also verified Lemma~\ref{lem:earlier} and the inequalities~\eqref{eq:two-sided} and~\eqref{eq:envelope}. The checks cover $1{,}111$ coefficient vectors and $2{,}446{,}000$ indices in total:
\begin{enumerate}
\item[(1)] all vectors with $k=2,3,4,5,6$ and entries at most $15,6,4,3,2$, respectively, with $2{,}000$ terms each;
\item[(2)] $96$ pseudo-random vectors with $2\le k\le16$ and entries between $1$ and $10^6$, with $3{,}000$ terms each;
\item[(3)] eleven special vectors, namely $(1,1)$, $(1,2)$, $(2,1)$, $(1,3)$, $(1,4)$, $(1,10^6)$, $(10^6,1)$, $(1,1,10^6)$ and $(10^6,1,1)$, together with the vector of dimension $20$ with all entries $1$ and the vector of dimension $20$ whose entries are $1$ except for the last entry $10^6$; each was checked for $10{,}000$ terms, except $(1,4)$, which was checked for $50{,}000$ terms.
\end{enumerate}
All checks passed. Table~\ref{tab:ranges} lists the observed ranges of $D_n$ for some vectors of dimension two. For $(1,1)$ the range agrees with Zaslavsky's formula, and for $(1,2)$ and $(2,1)$ with the bounds proved in~\cite[Theorems~1 and~2]{FJ}. The vectors $(1,3)$ and $(1,4)$ are not $A_1$-bounded, and for $(1,4)$ the positive value $D_{348}=1$ is the one discussed in the introduction.

\begin{table}[ht]
\caption{Observed ranges of $D_n=A_n-\kappa n$ for $1\le n\le N$.}\label{tab:ranges}
\begin{tabular}{lrrrcrr}
\toprule
$\ba$ & $A_1$ & $\kappa$ & $N$ & $A_1$-bounded & $\min D_n$ & $\max D_n$\\
\midrule
$(1,1)$ & $3$ & $5$ & $10{,}000$ & yes & $-2$ & $-1$\\
$(1,2)$ & $5$ & $7$ & $10{,}000$ & yes & $-3$ & $-1$\\
$(2,1)$ & $4$ & $7$ & $10{,}000$ & yes & $-4$ & $-1$\\
$(1,3)$ & $7$ & $9$ & $10{,}000$ & no & $-5$ & $0$\\
$(1,4)$ & $9$ & $11$ & $50{,}000$ & no & $-6$ & $1$\\
\bottomrule
\end{tabular}
\end{table}

\section*{Use of AI-assisted tools}

AI-assisted tools were used in the research and preparation of this manuscript.

\begin{thebibliography}{9}

\bibitem{B}
W.~Bosma, R.~Bruin, R.~Fokkink, J.~Grube, A.~Reuijl, and T.~Tromp, \emph{Using \texttt{Walnut} to solve problems from the OEIS}, J. Integer Seq. \textbf{28} (2025), Article 25.3.8.

\bibitem{FJ}
R.~Fokkink and G.~Joshi, \emph{Anti-recurrence sequences}, Integers \textbf{26} (2026), Paper No.~A24.

\bibitem{FP}
A.~S.~Fraenkel and U.~Peled, \emph{Harnessing the unwieldy MEX function}, Games of No Chance~4, Math. Sci. Res. Inst. Publ., vol.~63, Cambridge Univ. Press, 2015, 77--94.

\bibitem{KM}
C.~Kimberling and P.~J.~C.~Moses, \emph{Linear complementary equations and systems}, Fibonacci Quart. \textbf{57} (2019), no.~5, 96--110.

\bibitem{OEIS}
OEIS Foundation Inc., \emph{The On-Line Encyclopedia of Integer Sequences}, \url{https://oeis.org}.

\bibitem{Z}
T.~Zaslavsky, \emph{Anti-Fibonacci numbers: a formula}, unpublished note, 2016, \url{https://oeis.org/A075326/a075326_2.pdf}.

\end{thebibliography}

\end{document}
